\section{Proposed Research}

In this section, we will ...

\subsection{}

\subsection{Overview of ML Components for Cache and Memory Management}

% Here memory hierarchy is not the whole hierarchy
% TODO Move to intro
%The project will investigate the use of distributed machine leaning components
%from the last-level cache to the main memory, including hardware cache
%prefetching.  Those layers together play a major role in memory system
%performance and energy efficiency. They are also challenging to manage given
%their ever increasing complexity and the diversity of application demands.
%Traditional techniques in cache replacement, cache prefetching, and memory
%access scheduling have been extensively studied, but their gains are increasing
%diminishing.  The recent advance in ML techniques represents a new
%and promising approach to further improve cache and memory performance.

%However, there are unique challenges in applying techniques to optimize the
%cache and memory systems, and more broadly any components inside or partially
%inside a processor.  The ML component must be lightweight and cost-effective;
%otherwise, the performance gain can hardly be justified.  The research on
%performance-oriented NN accelerator, which is the main stream, does not have
%this constraint.  On those accelerator chips, the majority of the chip area is
%dedicated to NN acceleration.  By contrast, a processor chip should only budget
%a small percentage of its chip area for any ML circuits for improving cache and
%memory performance.  Otherwise, the chip area could be used to increase the
%cache capacity.  Additionally, the ML circuits must be energy-efficient.
%High-performance ML accelerator designs, by contrast, tend to be power hungry.

% For large L3 cache, should there be multiple ML components for the
% cache controller?

%Extensive research has been conducted on conventional hardware techniques for
%cache and memory management, including cache organization, multi-level cache
%design, cache replacement policies, cache prefetching, and main memory access
%scheduling.  Those efforts yielded significant results; however, any further
%effort is likely to have diminishing returns.  On the other hand, memory
%performance and memory energy efficiency continue to be a major hurdle for many
%memory-intensive workloads, even with the tremendous advances in cache capacity
%and memory data rates.  They must be complemented by more intelligent management
%techniques.  In fact, the increasing complexity of cache and memory also demands
%further advances in their management.  ML technique is promising in this regard
%because it can explore broader feature spaces for performance correlation.  It
%has been used to improve branch prediction~\cite{} and cache prefetching,
%showing the potential of ML technique to further improve performance on top of
%conventional techniques.  However, the use of ML is limited in scope, and the ML
%circuit design is still primitive.

% ML components should be added to L2 cache controller of each processor core.
% Assume that the processor has 3 levels of cache.

% Cache replacement, cache prefetching, memory access scheduling
% In their decision logic.
% Cache replacement: Affect LRU chain.  It decides by how much a block is prompted.
% What happens to non-LRU design? Pseudo LRU is OK.
% What feedback? If a prompted block is hit, then that's a positive feedback.
%
% One problem is that the feedback unit is expensive, because the past decision
% must be memorized. We can use a general structure for sampled memorization.
% Similar to counting network used in Internet. Each decision is encoded into
% a binary string, and then selectively memorized by sampling. Thus, the resource
% spent on feedback will be fixed.
%
% Any inexpensive feedback unit design?
% Sampling? Bloom filter? Pre-generate positive feedback and negative feedback?
% Cache prefetching: filter prefetching request. Need a feedback unit.
% In the current design (Texas A&M): check if prefetched block is used,
% or rejected block is hit. We need new ideas to show this project's value.
% Memory access scheduling:
%

% Overview

This project will apply ML techniques broadly to the cache and main memory
hierarchy by incorporating distributed and coordinated ML components through the
hierarchy.  It consists of three types of ML components to improve the decision
logic, specifically for cache replacement, cache prefetching, and memory access
scheduling.  They are embedded into the corresponding hardware components,
namely cache controller, hardware prefetcher, and memory controller.  They are
also connected with each other using the existing connections among those
components, and are coordinated to improve the learning and inference of each
other.

\subsection{ML Component for Cache Replacement}

\subsection{ML Component for Cache Prefetching}

\subsection{ML Component for Memory Access Scheduling}

\subsection{Coordination Framework for ML Components}

\subsection{Hardware and Software Coordination}

\subsection{}